% oneside and openany: chapters start on the next page, with no blank pages in between.
% For a printed, double-sided report use twoside,openright instead (and widen the inner margin).
\documentclass[11pt,a4paper,oneside,openany,titlepage,
               headings=small,bibliography=totoc,numbers=noenddot]{scrreprt}
\usepackage[T1]{fontenc}
\usepackage[utf8]{inputenc}
\usepackage{lmodern}
\usepackage[a4paper,left=3cm,right=2.5cm,top=2.8cm,bottom=2.8cm]{geometry}
\usepackage{microtype}
\usepackage{amsmath}
\usepackage{graphicx}
\usepackage{booktabs}
\usepackage{tabularx}
\usepackage{pgfplots}
\pgfplotsset{compat=1.16}
\usepackage[font=small,labelfont=bf]{caption}
\usepackage{enumitem}
\setlist{nosep}
\usepackage[hidelinks]{hyperref}
\usepackage[capitalise,noabbrev]{cleveref}

% ---- Fill in once: used on the title page, in the headers and in the PDF info ----
\newcommand{\reporttitle}{Title of the Report}
\newcommand{\reportsubtitle}{A subtitle that says what the report covers}
\newcommand{\reportnumber}{TR-2026-01}
\newcommand{\reportversion}{1.0}      % keep in step with the last row of the revision history
\newcommand{\reportdate}{15 October 2026}
\newcommand{\reportstatus}{Final}     % Draft, For review, Final ...
\newcommand{\reportauthors}{First Author, Second Author}
\newcommand{\reportorg}{Example Institute}
\hypersetup{pdftitle={\reporttitle}, pdfauthor={\reportauthors}}

% ---- Running headers and footers: report number left, chapter right ----------
\usepackage[headsepline,footsepline]{scrlayer-scrpage}
\clearpairofpagestyles
\automark[chapter]{chapter}
\ihead{\reportnumber{} v\reportversion}
\ohead{\headmark}
\cfoot{\pagemark}
\renewcommand{\chapterpagestyle}{scrheadings}  % keep the header on chapter-opening pages too

\begin{document}

% ---- Title page -------------------------------------------------------------
\begin{titlepage}
  \raggedright
  {\large\sffamily \reportorg\par}
  \vspace*{2.5cm}
  \rule{\textwidth}{1.2pt}\par\vspace{0.8cm}
  {\sffamily\bfseries\Huge \reporttitle\par}
  \vspace{0.6cm}
  {\sffamily\Large \reportsubtitle\par}
  \vspace{0.8cm}
  \rule{\textwidth}{0.4pt}\par\vspace{0.8cm}
  {\large \reportauthors\par}
  \vfill
  % The block that identifies this exact document: someone holding a printout can tell
  % which version it is and ask for the right one.
  \begin{tabular}{@{}l@{\qquad}l@{}}
    \sffamily Report number & \reportnumber \\
    \sffamily Version       & \reportversion \\
    \sffamily Date          & \reportdate \\
    \sffamily Status        & \reportstatus \\
    \sffamily Distribution  & Public \\
  \end{tabular}
\end{titlepage}

\pagenumbering{roman}

% ---- Document control: revision history and contact -----------------------------
\chapter*{Document control}
\addcontentsline{toc}{chapter}{Document control}
\markboth{Document control}{Document control}
\section*{Revision history}
Add a row for every version that leaves the team, newest last. Say what changed, not just
that something did.

\medskip
\noindent\begin{tabularx}{\textwidth}{@{}l l l X@{}}
  \toprule
  Version & Date & Author & Changes \\
  \midrule
  0.1 & 2026-08-03 & F.\ Author & First draft of chapters 1 and 2. \\
  0.2 & 2026-09-01 & S.\ Author & Added results; revised the approach after review. \\
  0.9 & 2026-09-28 & F.\ Author & Executive summary, appendices; ready for review. \\
  1.0 & 2026-10-15 & F.\ Author & Review comments addressed; released. \\
  \bottomrule
\end{tabularx}
% Dates in ISO 8601 order (year-month-day) cannot be misread across countries \cite{iso8601}.

\section*{Contact}
First Author, \reportorg, \texttt{first.author@example.org}.

% ---- Executive summary --------------------------------------------------------
% \addchap makes an unnumbered chapter that is still listed in the contents.
\addchap{Executive summary}
Write this for the reader who will read nothing else, and write it last. Use one page,
in this order:
\begin{itemize}
  \item \textbf{Purpose.} Why the work was done and for whom.
  \item \textbf{What was done.} The approach in two sentences.
  \item \textbf{Findings.} The three or four results that matter, with numbers.
  \item \textbf{Recommendations.} What the reader should do now, and by when.
\end{itemize}
Avoid abbreviations, citations and cross-references: the summary must make sense when it is
copied into an email.

\tableofcontents
\listoffigures
\listoftables

\clearpage
\pagenumbering{arabic}

% ---- Chapters ---------------------------------------------------------------
% One file per chapter keeps main.tex short and lets several people edit in parallel.
\chapter{Introduction}
\label{ch:intro}

\section{Background and purpose}
Explain why this report exists: the project, the question it answers and who asked. Two or
three paragraphs are enough; the executive summary has already given the short version.

\section{Scope}
Say what the report covers and, just as important, what it does not. A reader who looks for
the cost analysis should learn here that it is in another document.

\section{Structure of this report}
\Cref{ch:approach} describes the approach and \cref{ch:results} the results.
\Cref{ch:conclusions} gives the conclusions and recommendations. Supporting material is in
the appendices, starting with \cref{app:data}.

\chapter{Approach}
\label{ch:approach}

\section{Method}
Describe what was done in enough detail for a colleague to repeat it. Work that follows a
plan can be shown as a schedule \cite{gantt1919}; equations are numbered as usual:
\begin{equation}
  \eta = \frac{P_{\mathrm{out}}}{P_{\mathrm{in}}} \times 100\,\% .
  \label{eq:efficiency}
\end{equation}

\section{Materials and tools}
List software versions, equipment and data sources. A table is the clearest form:

\begin{table}[ht]
  \centering
  \caption{Tools used in the project.}
  \label{tab:tools}
  \begin{tabular}{@{}lll@{}}
    \toprule
    Tool & Version & Purpose \\
    \midrule
    Analysis scripts & 2.3 & Data processing \\
    Test rig         & Rev.\ B & Measurements \\
    \LaTeX{} \cite{lamport94} & 2026 & This report \\
    \bottomrule
  \end{tabular}
\end{table}

\chapter{Results}
\label{ch:results}

\section{Main findings}
Give the results in the order of the questions in \cref{ch:intro}, one subsection each,
and say what each result means in the same paragraph. \Cref{fig:trend} shows a figure drawn
in \LaTeX{} from data, so no image file is needed; replace it with
\verb|\includegraphics| for your own.

\begin{figure}[ht]
  \centering
  \begin{tikzpicture}
    \begin{axis}[width=0.85\textwidth, height=6cm, xlabel={Week}, ylabel={Efficiency (\%)},
                 xmin=0, xmax=12, ymin=60, ymax=100, grid=major, grid style={gray!25},
                 legend pos=south east]
      \addplot[thick, mark=*] coordinates {(1,68) (3,74) (5,81) (7,85) (9,88) (11,90)};
      \addplot[thick, dashed, domain=0:12] {85};
      \legend{Measured, Target}
    \end{axis}
  \end{tikzpicture}
  \caption{Efficiency $\eta$ from \cref{eq:efficiency} over the test period.}
  \label{fig:trend}
\end{figure}

\section{Comparison with the plan}
A table of planned against actual values lets the reader see at once what slipped.

\begin{table}[ht]
  \centering
  \caption{Planned and actual results.}
  \label{tab:plan}
  \begin{tabular}{@{}lrrl@{}}
    \toprule
    Item & Planned & Actual & Status \\
    \midrule
    Efficiency (\%) & 85 & 90 & Met \\
    Cost (k\texteuro) & 120 & 131 & Over budget \\
    Duration (weeks) & 12 & 11 & Met \\
    \bottomrule
  \end{tabular}
\end{table}

\chapter{Conclusions and recommendations}
\label{ch:conclusions}

\section{Conclusions}
Restate the findings that support the recommendations, without new results.

\section{Recommendations}
Number them, so that readers can refer to them in a meeting, and give each an owner and a date:
\begin{enumerate}
  \item Adopt the new process in all projects (owner: head of engineering, by 2027-01-31).
  \item Repeat the measurement with the second supplier (owner: project lead, by 2026-12-15).
\end{enumerate}

\section{Open issues}
Anything that is still unclear or unfinished, so that the next person does not rediscover it.


% ---- Appendices -----------------------------------------------------------------
\appendix
\chapter{Measurement data}
\label{app:data}
Raw data, long tables and derivations go in an appendix, so that the chapters stay readable.
Each appendix is a chapter, numbered A, B, and so on.

\begin{table}[ht]
  \centering
  \caption{Raw measurements.}
  \label{tab:raw}
  \begin{tabular}{@{}rrr@{}}
    \toprule
    Week & Input (W) & Output (W) \\
    \midrule
    1  & 100.0 & 68.2 \\
    5  & 100.0 & 81.0 \\
    11 & 100.0 & 90.1 \\
    \bottomrule
  \end{tabular}
\end{table}

\chapter{Glossary}
\label{app:glossary}
\begin{description}[style=nextline]
  \item[Efficiency] Output power divided by input power, in percent (\cref{eq:efficiency}).
  \item[Revision history] The table of versions at the start of this report.
\end{description}


% ---- References -----------------------------------------------------------------
\bibliographystyle{plain}
\bibliography{refs}

\end{document}
